\documentclass{standalone}
\usepackage{pgfplots}
\pgfplotsset{compat=1.18} % Assure la compatibilité avec la dernière version de pgfplots
\usepackage[utf8]{inputenc}
\usepackage[french]{babel} % Optionnel : pour la langue française
\usepackage{siunitx}

\usetikzlibrary{patterns}
\usetikzlibrary{arrows.meta}
\begin{document}

\begin{tikzpicture}
    \begin{axis}[
        xlabel={\( v_e \)},
        ylabel={Force},
        xmin=0, xmax=3000,
        ymin=0, ymax=7,
        xtick=\empty,
        ytick=\empty,
        scaled x ticks=false,
        tick label style={font=\scriptsize},
        ]

    % Courbe pour phi écranté
    \addplot[
        domain=0:3000,
        color=black,
        thick
    ]
    table [header=false] {force_Ee/force.dat};
    \node at (1700, 5) {\small \shortstack{\( mv_e/\tau_s \) Force\\collisionnelle}};
    
    % Courbe pour phi sans écrantage
    \addplot[
        domain=0:1000,
        color=black!70,
        thick
    ]
    table [header=false] {force_Ee/Ee.dat};
    \node at (2900, 1.8) {\( Ee \)};
    
    \def\xintercep{1935}
    \def\yintercep{2}
    \draw[ dashed ] (\xintercep, \yintercep-0.5) -- (\xintercep, \yintercep+1);
    \draw [ -{Latex} ] (\xintercep, \yintercep+0.5) -- (\xintercep+700, \yintercep+0.5) node[ midway, anchor=south ] {\small Fugitif};
    \fill (\xintercep, \yintercep) circle  [ radius=0.075cm ];
    \end{axis}
\end{tikzpicture}

\end{document}
